% IEEE conference version for GTSD 2026 proceedings
\documentclass[conference]{IEEEtran}
\usepackage{cite}
\usepackage{amsmath,amssymb,amsfonts}
\usepackage{graphicx}
\graphicspath{{figures/}}
\usepackage{textcomp}
\usepackage{url}
\def\BibTeX{{\rm B\kern-.05em{\sc i\kern-.025em b}\kern-.08em
    T\kern-.1667em\lower.7ex\hbox{E}\kern-.125emX}}

\begin{document}
\bstctlcite{IEEEexample:BSTcontrol}

\title{An Energy-Efficient and Resource-Constrained FPGA Accelerator for Medical Image Super-Resolution Using Quantized SRCNN}

\author{\IEEEauthorblockN{Le Van Duc}
\IEEEauthorblockA{\textit{Faculty of Information Technology}\\
\textit{VNU-HCM University of Science}\\
Ho Chi Minh City, Viet Nam\\
24280010@student.hcmus.edu.vn}
}
\maketitle
\section{Introduction}

Medical imaging modalities, including X-ray radiography, Computed Tomography (CT), and Magnetic Resonance Imaging (MRI), serve as indispensable diagnostic cornerstones in modern healthcare. The clinical efficacy of these imaging systems relies heavily on image quality, particularly spatial resolution, contrast, noise characteristics, and artifact suppression, which collectively determine the visibility and interpretability of anatomical and pathological structures \cite{obuchowicz2026image, palma2021review}. For instance, the ability to resolve fine anatomical details and subtle lesions can be critical for accurate detection and characterization of pathology \cite{obuchowicz2026image, palma2021review}. However, acquiring high-quality HR images directly during clinical routines is constrained by acquisition time, patient motion, radiation exposure, and the physical limitations of imaging systems \cite{chen2023deep, shin2024super, occhipinti2026ultra}. In MRI, increasing spatial resolution can require longer acquisition times, lower signal-to-noise ratios, and can increase susceptibility to motion artifacts \cite{chen2023deep}. For X-ray and CT, radiation exposure imposes an additional constraint, motivating the use of low-dose acquisition strategies that may introduce increased noise and image artifacts \cite{shin2024super}. Furthermore, the achievable spatial resolution of CT is fundamentally constrained by hardware and acquisition factors such as detector element size, focal spot size, projection sampling, and acquisition geometry \cite{occhipinti2026ultra}. Upgrading imaging hardware can therefore require substantial capital investment while remaining subject to these physical constraints.

To bypass these hardware restrictions, Single Image Super-Resolution (SISR) has emerged as an effective post-acquisition computational paradigm, aiming to recover lost high-frequency spatial components and reconstruct high-fidelity HR representations from Low-Resolution (LR) observations \cite{yang2019deep, liu2025comprehensive}. Nevertheless, applying SR to medical modalities presents profound technical challenges. Fundamentally, SISR is an ill-posed inverse problem in which a single LR observation can correspond to multiple plausible HR solutions \cite{yang2019deep, liu2025comprehensive}. Furthermore, medical images can exhibit limited soft-tissue contrast and reduced signal-to-noise ratios depending on the imaging modality and acquisition protocol, making unconstrained reconstruction susceptible to noise amplification and artifact generation \cite{obuchowicz2026image, varghese2024enhancing}. Consequently, designing algorithms that can robustly enhance structural discontinuities and preserve fine diagnostic edges without generating artificial anatomical hallucinations remains a critical necessity. In particular, recent studies on medical image reconstruction and MRI super-resolution have highlighted the risk that deep learning models may introduce anatomically implausible structures or hallucinated high-frequency features, potentially compromising diagnostic reliability \cite{bhadra2021hallucinations, giraldo2024perceptual}.

Driven by the rapid evolution of deep learning, learning-based super-resolution frameworks have achieved remarkable improvements in reconstruction fidelity across widely used benchmark evaluations \cite{dong2014learning, kim2016accurate, lim2017enhanced}. In the literature, deep SISR architectures have evolved through several major paradigms. Early pioneering methodologies utilized compact Convolutional Neural Networks (CNNs) to learn non-linear mappings directly from low-resolution inputs to high-resolution outputs, exemplified by SRCNN \cite{dong2014learning}, FSRCNN \cite{dong2016accelerating}, and ESPCN \cite{shi2016real}. To improve representation capability and facilitate the training of deeper architectures, residual learning was subsequently introduced in networks such as VDSR \cite{kim2016accurate} and EDSR \cite{lim2017enhanced}, substantially improving pixel-level reconstruction fidelity. More recently, Generative Adversarial Networks (GANs), exemplified by SRGAN \cite{ledig2017photo}, have been investigated to recover perceptually realistic high-frequency textures through adversarial and perceptual objectives. In parallel, Vision Transformer-based architectures, such as SwinIR \cite{liang2021swinir}, have demonstrated strong performance for image super-resolution by exploiting long-range feature interactions through hierarchical self-attention.

Despite their impressive reconstruction performance, many high-performing SR models introduce substantial computational, parameter, and memory requirements, creating significant challenges for deployment on resource-constrained platforms. This issue is particularly pronounced for embedded and edge devices, where limited computational capability, memory capacity, and power budgets constrain the practical deployment of deep SR networks \cite{wang2023lightweight, xiao2022mfen}. Recent studies on on-device SR further demonstrate that high-resolution inference can impose considerable latency and computational costs, motivating the development of parameter-efficient and hardware-aware architectures \cite{hadji2025edge}. These constraints become even more critical for FPGA-based implementations, where computational complexity, operator efficiency, and memory bandwidth directly affect achievable throughput and hardware resource utilization \cite{huang2026uasr}. Consequently, bridging the gap between high-fidelity SR reconstruction and efficient edge inference remains an important research challenge.

To overcome this problem, many researchers have proposed various architectures and techniques to accelerate the inference of CNNs \cite{huynh2022fpga}. With respect to hardware implementations, three hardware platforms can be used as CNN accelerators: graphic processing units (GPU) \cite{strigl2010performance}, application-specific integrated circuits (ASICs) \cite{chen2016eyeriss}, and field programmable gate arrays (FPGAs) \cite{zhang2015optimizing, sit2017fpga}. Among these platforms, FPGA has revealed itself as the high-performance and low-cost embedded device very suitable for the hardware prototyping of convolutional accelerators. Recently, new FPGA hardware solutions have been introduced, allowing for an efficient hardware and software co-design framework for deep learning applications. The PYNQ \cite{pynqhomepage} is an open-source project from Xilinx that makes it easier to develop FPGA-based deep learning applications, in which designers can efficiently combine the benefits of programmable logic and microprocessors using the Python language and libraries.

Motivated by these challenges, this paper presents an end-to-end hardware-software co-design framework for accelerating medical image super-resolution on a resource-constrained edge FPGA platform. Specifically, we propose a lightweight 3-layer SRCNN architecture tailored for low-power edge deployment on the Xilinx Zynq-7000 SoC (PYNQ-Z2). To eliminate floating-point arithmetic overhead and comply with the physical constraints of only 220 DSP slices on the target XC7Z020 chip, we apply Quantization-Aware Training (QAT) to compress the model into a signed 8-bit fixed-point (Q7) representation. On the hardware side, an optimized streaming RTL accelerator is designed with line-buffers, parallel MAC units, and an adder tree, seamlessly interfaced with the host ARM processing system via high-throughput AXI4-Stream DMA transfers. Furthermore, an overlap-and-save tiling mechanism is integrated on the host to eliminate boundary blocking artifacts on large medical radiographs. The main contributions of this work are summarized as follows:
\begin{itemize}
    \item We propose a hardware-friendly, lightweight 3-layer SRCNN architecture and establish a Quantization-Aware Training (QAT) Q7 scheme that effectively minimizes quantization errors and preserves critical anatomical details in medical images.
    \item We design a fully pipelined RTL convolutional accelerator on the Xilinx Zynq-7000 (XC7Z020) FPGA, achieving high hardware compute efficiency by utilizing on-chip line buffers and adder trees while strictly respecting the 220 DSP budget.
    \item We implement an end-to-end PYNQ-based deployment with AXI DMA streaming and overlap-and-save tiling, validating the system across comprehensive clinical benchmarks (NIH ChestX-ray14 and Pediatric Pneumonia datasets) in terms of reconstruction fidelity, inference latency, and energy efficiency.
\end{itemize}

The remainder of this paper is organized as follows. Section~\ref{sec:related_work} reviews the related work on medical image super-resolution algorithms and FPGA-based deep learning accelerators. Section~\ref{sec:proposed_system} details the proposed hardware-software co-design framework, including the QAT Q7 quantization scheme and the RTL accelerator architecture. Section~\ref{sec:result_analysis} presents the experimental results and performance evaluation across objective image quality metrics, hardware resource utilization, and energy efficiency. Finally, Section~\ref{sec:conclusion} concludes the paper and outlines directions for future research.

\bibliographystyle{IEEEtran}
\bibliography{bib/references}

\end{document}